\section{Defects propagate as a charge-holonomy automaton}

\Cref{thm:defect-is-inconsistency-sheaf} supplies the defect fiber: a section of
the $\top$-sheaf over the interlacement order.  This subgraph adds the missing
transition carrier: charge holonomy, cascade propagation, and the circuit reading
of controlled defects.

\begin{theorem}[The disclination is a holonomy quantum]\label{thm:disclination-holonomy-quantum}\lean{Obsmetrologia.AperiodicOrder.Defect.Disclination.disclination_holonomy_quantum}\leanok
On the hexagonal carrier — frames as vectors of the Eisenstein-integer plane turned
by the derived sixth-rotation, whose orbit has order exactly six as a theorem, not a
declared carrier — a six-coordinated site's boundary loop develops closed (flat),
while a five-coordinated site's loop returns the frame rotated: one quantum of
rotational holonomy, and the development first closes at the sixth traversal. A
seven-site carries the opposite charge: a 5--7 pair at zero separation annihilates
(the joint circuit develops to the identity), while a 5--7 pair at positive
separation is a dislocation — trivial rotational holonomy, nonzero translational
holonomy, with Burgers vector growing linearly in the separation
(\texttt{dipole\_dislocation}: $(0,-1), (0,-2), (0,-3)$ at separations one, two,
three). Rotational holonomy is curvature and translational holonomy is torsion (the
crystal frontier memo, \S 5.5); the disclination is the discrete curvature quantum,
the dislocation its confined dipole. The development depends on no axioms.
\end{theorem}

\begin{theorem}[The Burgers vector is linear in the separation]\label{thm:dipole-burgers-linear}\lean{Obsmetrologia.AperiodicOrder.Defect.Disclination.dipole_burgers_linear}\leanok\uses{thm:disclination-holonomy-quantum}
At every separation, the 5--7 dipole circuit develops to the identity frame
with the position displaced by exactly $(0, -\mathrm{sep})$
(\texttt{dipole\_burgers\_linear}), and the Burgers vector vanishes iff the
separation is zero (\texttt{dipole\_burgers\_zero\_iff\_contact}): the
pair's dislocation content is the separation, uniquely minimized at
contact, where the pair annihilates. Proved by translation equivariance of
the development (\texttt{develop\_shift}) and the straight-walk law
(\texttt{develop\_straight}); the instance values of
\cref{thm:disclination-holonomy-quantum} are the cases one, two, three.
\end{theorem}

\begin{theorem}[The obstruction functional is additive]\label{thm:config-burgers-additive}\lean{Obsmetrologia.AperiodicOrder.Defect.Disclination.config_burgers_additive}\leanok\uses{thm:dipole-burgers-linear}
A configuration of dipoles — the concatenated circuit of one 5--7 dipole
per listed separation — develops to the identity frame with the position
displaced by the sum of the separations
(\texttt{config\_burgers\_additive}): the many-body dislocation content is
the sum of the two-body contents, and it vanishes iff every separation is
zero (\texttt{config\_burgers\_zero\_iff\_contact}) --- each pair
annihilated, none merely cancelled against another. The many-body
potential of \cref{thm:gravity-obstruction-descent} is additive on this
carrier.
\end{theorem}

\begin{theorem}[Obstruction descent settles at the vacuum]\label{thm:obstruction-descent-dynamics}\lean{Obsmetrologia.AperiodicOrder.Defect.Disclination.obstruction_descent}\leanok\uses{thm:config-burgers-additive}
The slide dynamics on dipole configurations --- every pair moves one unit
toward contact per stage. The total obstruction is non-increasing along
the flow (\texttt{sum\_slide\_le}), strictly decreasing at every
configuration with an unannihilated pair (\texttt{sum\_slide\_lt}), its
fixed points are exactly the annihilated configurations
(\texttt{slide\_fixed\_iff\_vacuum}), and the descent reaches zero total
Burgers vector in as many stages as the largest separation
(\texttt{slide\_collapse}): the settling stage of the collapse is the
configuration's diameter. The dynamics clauses depend on no axioms.
\end{theorem}

\begin{theorem}[The rotational charge is conserved along the descent]\label{thm:disclination-charge-conservation}\lean{Obsmetrologia.AperiodicOrder.Defect.Disclination.unpaired_disclination_survives_descent}\leanok\uses{thm:obstruction-descent-dynamics, thm:disclination-holonomy-quantum}
The developed frame of any walk is the iterated rotation by the walk's
total turn count (\texttt{develop\_fst\_eq\_rotPow\_sum}), composing along
concatenation (\texttt{develop\_fst\_append}), and straight segments cannot
move it. The turn counts act through the frame orbit, on which the derived
rotation has order six (\texttt{rotPow\_six\_of\_frameMem}, the orbit-wide
form of \texttt{frame\_orbit\_six}): a based walk's frame closes exactly
when its turn count is a multiple of six
(\texttt{develop\_frame\_closed\_iff}), a $k$-site's boundary loop is
frame-flat exactly when its coordination $k$ is
(\texttt{cone\_frame\_closed\_iff} --- the coordination census), and a
configuration of cones is frame-flat exactly when its total coordination
is (\texttt{coneConfig\_frame\_closed\_iff} --- the neutral sectors; the
frame component of \texttt{five\_seven\_annihilate} is the $[5,7]$
instance). A dipole circuit is frame-neutral at every separation because
its charge word cancels --- the five-cone's deficit quantum against the
seven-cone's excess quantum (\texttt{rotPow\_seven\_of\_frameMem}), the two
half-turns against each other (\texttt{rotPow\_uturn\_pair\_of\_frameMem})
--- so the slide dynamics, which edits only straight segments, conserves
the rotational charge along the whole descent: a configuration containing
one unpaired five-site develops to the same one-quantum frame at every
slide stage (\texttt{unpaired\_frame\_invariant}) and never closes. The
capstone bundle: the descent discharges the configuration's entire
translational content while the unpaired rotational quantum is fixed at
every stage --- the settled value of the collapse carries a conserved
remainder, and only an opposite cone annihilates it, at any separation
(\texttt{seven\_neutralizes\_unpaired}). All clauses except the inherited
\texttt{slide\_collapse} component depend on no axioms.
\end{theorem}

\begin{theorem}[Attraction is obstruction descent]\label{thm:gravity-obstruction-descent}\notready\uses{thm:dipole-burgers-linear, thm:config-burgers-additive, thm:obstruction-descent-dynamics, thm:anti-landauer-index-balance, thm:defect-is-inconsistency-sheaf}
\emph{(Open frontier --- stated, not formalized.)} The gravity reading of
the defect calculus. A mass is an entrenched local section of the
inconsistency sheaf (\cref{thm:defect-is-inconsistency-sheaf}) that the KZ
closure cannot flatten --- a conflict site no local substitution removes.
An obstruction functional on defect configurations (the \v{C}ech/Landauer
ledger of \cref{thm:anti-landauer-index-balance} summed over the boundary
link graphs of super-tiles, the Einstein--Hilbert analog of the M\"obius
incidence identity \texttt{total\_ricciScalar\_eq\_two\_bdAction}) is
lowered by moving conflict sites together: by
\cref{thm:dipole-burgers-linear} the pair's dislocation content is linear
in the separation and vanishes at contact, so any obstruction-decreasing
(anti-Landauer descent) dynamics drives defects toward mutual
annihilation. Attraction is the descent direction of the obstruction
ledger. The multi-dipole obstruction functional is supplied by
\cref{thm:config-burgers-additive} and the descent dynamics with its
Lyapunov law by \cref{thm:obstruction-descent-dynamics}; what remains is
the weld --- identify the slide with a settling closure of the KZ family
on a genuine tiling carrier, an entrenched inconsistency-sheaf section
with a disclination, and the obstruction sum with the boundary-link
\v{C}ech/Landauer ledger --- and the continuum reading of
\texttt{total\_ricciScalar\_eq\_two\_bdAction}.
\end{theorem}

\begin{theorem}[Defect unbinding is two-stage melting]\label{thm:kthny-defect-unbinding}\notready\uses{thm:disclination-holonomy-quantum, thm:real-ratio-temperature-connectivity-frontier}
\emph{(Open frontier --- stated, not formalized.)} The KTHNY ladder on a finite
defect ensemble: below a first balanced-flux temperature all 5--7 pairs are bound
(dislocations confined; both translational and sixfold orientational order);
between the two temperatures dislocations unbind but disclinations stay paired
(translational order lost, sixfold frame order survives --- the hexatic phase);
above the second temperature free disclinations deconfine and the frame order
collapses to the trivial group. Two transitions, each the deconfinement of one
holonomy charge of \cref{thm:disclination-holonomy-quantum}; the four-corner
reading (both orders / orientational only / positional only / neither) has its
positional-only corner physically empty. Requires a stochastic ensemble over the
disclination carrier with the balanced-flux temperature of the thermodynamics
sector.
\end{theorem}

\begin{theorem}[A defect is a charge-holonomy automaton]\label{thm:defect-charge-holonomy-automaton}\notready\uses{thm:defect-is-inconsistency-sheaf, thm:mobius-band-holonomy, thm:disclination-holonomy-quantum}
\emph{(Open frontier --- stated, not formalized.)} The defect-propagation calculus
(Shimamoto's dropped-loop classifier) is a finite automaton on the $\mathbb{Z}^2$ carrier
whose state is the accumulated charge holonomy around a loop enclosing the defect: a
mis-crossing carries $\top_k$, a hole carries $\bot_k$, and the transition rule transports
this charge to the next cell along the interlacement order. A defect is removable iff its
holonomy is trivial, exactly as the M\"obius band's face-swap holonomy
(\cref{thm:mobius-band-holonomy}) is the obstruction that cannot be combed away. Requires a
charge-transport automaton over the torus carrier.
\end{theorem}

\begin{theorem}[Laddering is a charge cascade]\label{thm:laddering-charge-cascade}\notready\uses{thm:defect-charge-holonomy-automaton, thm:unweaving-is-landauer}
\emph{(Open frontier --- stated, not formalized.)} A dropped stitch that ladders is the
charge-holonomy automaton (\cref{thm:defect-charge-holonomy-automaton}) in a non-settling
run: each released support relation propagates a $\bot_k$ (told-nothing) cell down its wale,
and the cascade halts iff the local contact graph regains forest structure --- the
un-weaving Landauer threshold (\cref{thm:unweaving-is-landauer}). A run-stop is a
cycle-rich, high-$\check H^1$ region that arrests the ladder; a clean ladder is a near-forest
that erases freely. Laddering propagation is the dual of settling: information un-sticking
along the wale.
\end{theorem}

\begin{theorem}[Unraveling gates compute on the $\top$-mirror]\label{thm:unraveling-gates-paraconsistent}\notready\uses{thm:defect-charge-holonomy-automaton, thm:paraconsistentEq-eq}
\emph{(Open frontier --- stated, not formalized.)} Shimamoto's unraveling logic gates ---
NOT, AND, OR, half-adder implemented by controlled defect propagation on stitch topology ---
are the charge-holonomy automaton (\cref{thm:defect-charge-holonomy-automaton}) used as a
circuit: each gate is a local transport rule on the held $\top$-mirror, computing
paraconsistently (\cref{thm:paraconsistentEq-eq}) from the contradiction without explosion.
Requires the automaton of \cref{thm:defect-charge-holonomy-automaton} read as a
paraconsistent circuit.
\end{theorem}
